\subsection{级数部分和的渐近展开}

对实变量 \(x\) （趋于 \(+\infty\) ），设 \(\mathcal{E}\) 是一个比较尺度，由这样的函数组成：其中每个函数都在整个区间 \([x_0, +\infty[\) （随函数而定）上有定义，且在该区间上 \(\geqslant 0\) 。设 \((\mathbf{u}_n)\) 是各项都属于完备赋范空间 \(E\) 的级数，使得 \(\mathbf{u}_n\) 有精确到 \(g_\alpha\) 的渐近展开，它相对于尺度 \(\mathcal{E}'\) （由 \(\mathbf{N}\) 上的、来自 \(\mathcal{E}\) 中函数的限制组成）：

\[
\mathbf {u} _ {n} = \sum_ {\lambda \leqslant \alpha} \mathbf {a} _ {\lambda} g _ {\lambda} (n) + \mathbf {r} _ {\alpha} (n).
\]

假设每个部分和 \(\sum_{m=1}^{n} g(m)\) （其中 \(g \in \mathcal{E}\) ）都有关于 \(\mathcal{E}'\) 的渐近展开。于是可以得到 \(\mathbf{s}_n = \sum_{m=1}^{n} \mathbf{u}_m\) 关于 \(\mathcal{E}'\) 的渐近展开；我们仍然区分两种情形：

\(1^{\circ}\sum_{n = 1}^{\infty}g_{\alpha}(n) = +\infty .\) 这时 (V,p.237,prop.2) 有 \(\sum_{m = 1}^{n}\mathbf{r}_{\alpha}(m)\ll \sum_{m = 1}^{n}g_{\alpha}(m);\)

根据假设，可以得到

\[
\sum_ {\lambda \leqslant \alpha} \mathbf {a} _ {\lambda} \left(\sum_ {m = 1} ^ {n} g (m)\right)
\]

(V, p.~222) 到某个精确度 \(g_{\sigma}\) 的渐近展开；若 \(cg_{\sigma}(n)\) 是 \(\sum_{m=1}^{n} g_{\alpha}(m)\) 的主部，则将有 \(\mathbf{s}_n\) 的精确到 \(g_{\min(\rho, \sigma)}\) 的渐近展开。

\(2^{\circ}\sum_{n = 1}^{\infty}g_{\alpha}(n)\) 收敛；这时设 \(\beta\) 是指标 \(\lambda \leqslant \alpha\) 中最小的一个，它使 \(\mathbf{a}_{\lambda}\neq 0\) 且使 \(\sum_{n = 1}^{\infty}g_{\lambda}(n)\) 收敛；则级数

\[
\mathbf {C} = \sum_ {n = 1} ^ {\infty} \left(\mathbf {u} _ {n} - \sum_ {\lambda <   \beta} \mathbf {a} _ {\lambda} g _ {\lambda} (n)\right)
\]

绝对收敛，并且可以写

\[
\mathbf {s} _ {n} = \sum_ {\lambda <   \beta} \mathbf {a} _ {\lambda} \left(\sum_ {m = 1} ^ {n} g _ {\lambda} (m)\right) + \mathbf {C} - \sum_ {\beta \leqslant \lambda \leqslant \alpha} \mathbf {a} _ {\lambda} \left(\sum_ {m = n + 1} ^ {\infty} g _ {\lambda} (m)\right) - \sum_ {m = n + 1} ^ {\infty} \mathbf {r} _ {\alpha} (m).
\]

此外，有 \(\sum_{m=n+1}^{\infty}\mathbf{r}_{\alpha}(m)\ll\sum_{m=n+1}^{\infty}g_{\alpha}(m)\) ；若 \(cg_{\sigma}(n)\) 是 \(\sum_{m=n+1}^{\infty}g_{\alpha}(m)\) 的主部，且若已有

\[
\sum_ {\lambda <   \beta} \mathbf {a} _ {\lambda} \left(\sum_ {m = 1} ^ {n} g _ {\lambda} (m)\right) + \mathbf {C} - \sum_ {\beta \leqslant \lambda \leqslant \alpha} \mathbf {a} _ {\lambda} \left(\sum_ {m = n + 1} ^ {\infty} g _ {\lambda} (m)\right)
\]

到精确度 \(g_{\rho}\) 的渐近展开，便由此得到 \(s_{n}\) 的精确到 \(g_{\min(\rho,\sigma)}\) 的渐近展开。于是问题归结为级数 \((g(n))\) （其中 \(g \in E\) ）的特殊情形。我们将看到，在某些条件之下，如何能直接得到 \(s_{n} = \sum_{m=1}^{n} g(m)\) 的主部（当 \(\sum_{n=1}^{\infty} g(n) = +\infty\) 时）或 \(r_{n} = \sum_{m=n+1}^{\infty} g(m)\) 的主部（当 \(\sum_{n=1}^{n} g(n) < +\infty\) 时）。

命题 6. 设 \(g\) 是一个实函数， \(>0\) 且单调，定义在区间 \([x_0, +\infty[\) 上（其中 \(x_0 \leqslant 1\) ），且使 \(\log g\) 与 \(x\) 1 阶可比较。

\(1^{\circ}\) 若 \(g\) 相对于 \(e^{x}\) 是无穷阶的，则有

\[
s _ {n} = \sum_ {m = 1} ^ {n} g (m) \sim g (n) \quad i f \quad \sum_ {n = 1} ^ {\infty} g (n) = + \infty ;\tag{1}
\]

\[
r _ {n} = \sum_ {m = n + 1} ^ {\infty} g (m) \sim g (n + 1) \quad i f \quad \sum_ {n = 1} ^ {\infty} g (n) <   + \infty .\tag{2}
\]

\(2^{\circ}\) 若 \(g\) 有有限阶 \(\mu\) （相对于 \(e^{x}\) ），则有

\[
s _ {n} = \sum_ {m = 1} ^ {n} g (m) \sim \frac {\mu}{1 - e ^ {- \mu}} \int_ {x _ {0}} ^ {n} g (t) d t \quad i f \quad \sum_ {n = 1} ^ {\infty} g (n) = + \infty ;\tag{3}
\]

\[
r _ {n} = \sum_ {m = n + 1} ^ {\infty} g (m) \sim \frac {\mu}{1 - e ^ {- \mu}} \int_ {n} ^ {\infty} g (t) d t \quad i f \quad \sum_ {n = 1} ^ {\infty} g (n) <   + \infty\tag{4}
\]

（数 \(\frac{\mu}{1 - e^{-\mu}}\) 在 (3) 与 (4) 中当 \(\mu = 0\) 时应换成 1）。

\(1^{\circ}\) 若 \(g\) 的阶为 \(+\infty\) （相对于 \(e^{x}\) ），则有 \(\log g \gg x\) ，从而按假设有 \(g' / g \gg 1\) ，即 \(g' \gg g\) ；因此 \(g\) 递增并趋于 \(+\infty\) （随 \(x\) ），从而 \(\sum_{n=1}^{\infty} g(n) = +\infty\) 。若 \(u\) 是与级数 \((g(n))\) 相关联的阶梯函数 (V, p.~237)，则 \(u(x) \leqslant g(x)\) 从 \(x\) 的某个值起成立，故 \(u \preccurlyeq g\) ，因而

\[
s _ {n - 1} = \int_ {1} ^ {n} u (t) d t \preccurlyeq \int_ {1} ^ {n} g (t) d t \prec \prec \int_ {1} ^ {n} g ^ {\prime} (t) d t \sim g (n);
\]

由于 \(s_n = s_{n-1} + g(n)\) ，有 \(s_n \sim g(n)\) 。当 \(g\) 的阶为 \(-\infty\) （相对于 \(e^x\) ）时，证明类似；这样便得到公式 (2)。

\(2^{\circ}\) 若 \(g\) 有有限阶 \(\mu\) （相对于 \(e^{x}\) ），则可写 \(g(x) = e^{\mu x}h(x)\) ，其中 \(h\) 相对于 \(e^{x}\) 是 0 阶的；此外，按假设， \(\log g \sim \mu x\) （对 \(\mu \neq 0\) ； \(\log g \prec x\) 对 \(\mu = 0\) ）蕴含 \(h' \prec h\) 。先设 \(\sum_{n=1}^{\infty} g(n) = +\infty\) （这蕴含 \(\mu \geqslant 0\) ；当 \(\mu > 0\) 时逆命题总成立，因为这时 \(g(x)\) 趋于 \(+\infty\) （随 \(x\) ））；我们来求 \(\int_{n-1}^{n} g(t) dt\) 的主部。可以写

\[
\begin{array}{l} \int_ {n - 1} ^ {n} g (t) d t = \int_ {n - 1} ^ {n} e ^ {\mu t} h (t) d t \\ \qquad = h (n) \int_ {n - 1} ^ {n} e ^ {\mu t} d t + \int_ {n - 1} ^ {n} e ^ {\mu t} \big (h (t) - h (n) \big) d t \\ \qquad = \frac {1 - e ^ {- \mu}}{\mu} g (n) + \int_ {n - 1} ^ {n} e ^ {\mu t} \big (h (t) - h (n) \big) d t. \end{array}
\]

现在，关系 \(h^{\prime} \ll h\) 蕴含：对每个 \(\varepsilon > 0\) ，存在 \(n_{0}\) ，使关系 \(x \geqslant n_{0}\) 蕴含 \(\left|h'(x)/h(x)\right| \leqslant \varepsilon\) ；由中值定理推出 \(-\varepsilon \leqslant \log |h(t)/h(n)| \leqslant \varepsilon\) （对 \(n - 1 \leqslant t \leqslant n\) ，当 \(n \geqslant n_{0}\) 时），由此

\[
| h (t) - h (n) | \leqslant \left(e ^ {\varepsilon} - 1\right) h (n)
\]

从而

\[
\left| \int_ {n - 1} ^ {n} e ^ {\mu t} (h (t) - h (n)) d t \right| \leqslant (e ^ {\varepsilon} - 1) e ^ {\mu n} h (n) = (e ^ {\varepsilon} - 1) g (n)
\]

因为 \(e^{\mu t}\) 递增。由于 \(e^{\varepsilon} - 1\) 随 \(\varepsilon\) 变得任意小，可见可以写

\[
\int_ {n - 1} ^ {n} g (t) d t = \frac {1 - e ^ {- \mu}}{\mu} g (n) + o (g (n))
\]

（\(\left(\frac{1 - e^{-\mu}}{\mu}\right.\) 当 \(\mu = 0\) 时换成 1。于是该命题是 V, p.~237 的命题 2 的推论。当 \(\sum_{n=1}^{\infty} g(n)\) 有限时，论证类似。

反复应用 V, p.~239 的命题 6，有时便能得到 \(s_{n} = \sum_{m=1}^{n} g(m)\) 的渐近展开。首先设 g 的阶为 \(+\infty\) （相对于 \(e^{x}\) ）；对 p 的每个固定值，由命题 6 可以写

\[
s _ {n} = g (n) + g (n - 1) + \dots + g (n - p) + o (g (n - p))
\]

只需（关于 \(E'\) ）展开每个函数 \(g(n-k)\) （ \(0 \leqslant k \leqslant p\) ），并把展开的精确度限制到 \(g(n-p)\) 的主部，便得到 \(s_{n}\) 的展开。

例. 设 \(g(x) = x^{x} = \exp (x\log x)\) ，其阶为 \(+\infty\) （相对于 \(e^x\) ）。取 \(p = 2\) ，有

\[
(n - 1) \log (n - 1) = (n - 1) \log n - 1 + \frac {1}{2 n} + o \left(\frac {1}{n}\right)
\]

由此 (V, p.~225)

\[
(n - 1) ^ {n - 1} = \frac {1}{e} n ^ {n - 1} + \frac {1}{2 e} n ^ {n - 2} + o _ {1} (n ^ {n - 2})
\]

同理

\[
(n - 2) ^ {n - 2} = \frac {1}{e ^ {2}} n ^ {n - 2} + o _ {2} (n ^ {n - 2});
\]

因此

\[
s _ {n} = n ^ {n} + \frac {1}{e} n ^ {n - 1} + \left(\frac {1}{2 e} + \frac {1}{e ^ {2}}\right) n ^ {n - 2} + o _ {3} (n ^ {n - 2}).
\]

对 \(r_n\) ，当 \(g\) 的阶为 \(-\infty\) （相对于 \(e^x\) ）时，同样处理。

现在若 \(g\) 有有限阶 \(\mu\) （相对于 \(e^x\) ），且例如 \(\sum_{n=1}^{\infty} g(n) = +\infty\) ，则可以写

\[
s _ {n} = \frac {\mu}{1 - e ^ {- \mu}} \int_ {1} ^ {n} g (t) d t + \sum_ {m = 1} ^ {n} f _ {1} (m)
\]

其中 \(f_{1}(n) = g(n) - \frac{\mu}{1 - e^{-\mu}}\int_{1}^{n}g(t)dt \ll g(n)\) （由 V, p.~239 的命题 6）。若已有 \(cg_{1}(n)\) 作为 \(f_{1}(n)\) 关于 \(\mathcal{E}'\) 的主部，且若能对函数 \(g_{1}\) 再应用命题 6，则将得到与 \(\sum_{m=1}^{n}f_{1}(m)\) 等价的原函数（当 \(\sum_{n=1}^{\infty}g_{1}(n) = +\infty\) 时），而在相反的情形得到与 \(\sum_{m=n+1}^{\infty}f_{1}(m)\) 等价的原函数（在后一情形写 \(\sum_{m=1}^{n}f_{1}(m) = C - \sum_{m=n+1}^{\infty}f_{1}(m)\) ，其中 \(C = \sum_{n=1}^{\infty}f_{1}(n)\) ）。

这样逐步进行，最终可以把 \(s_{n}\) 表示成若干个原函数之和，其中每个原函数相对于前一个是可忽略的，外加一项相对于已写出的最后一个原函数仍为可忽略的项，最后还有一个常数（余项趋于 0 的情形）。然后只需关于 \(E'\) 展开所得到的每个原函数 (cf.~V, p.~235)。

例. 设 \(g(n) = \frac{1}{n}\) ；这时

\[
s _ {n} = \sum_ {m = 1} ^ {n} \frac {1}{m} \sim \int_ {1} ^ {n} \frac {d t}{t} = \log n
\]

继而

\[
\frac {1}{n} - (\log n - \log (n - 1)) \sim - \frac {1}{2 n ^ {2}}
\]

由此

\[
s _ {n} = \log n + \gamma + \frac {1}{2 n} + o \left(\frac {1}{n}\right).
\]

出现在这个公式中的常数 \(\gamma\) 在分析中扮演着重要角色 (cf.~chap.~VI 与 VII)；它称为欧拉常数；有

\[
\gamma = 0. 5 7 7   2 1 5   6 6 4 \dots
\]

精确到 \(1/10^{9}\) 。

我们将在 VI, p.~288 看到，欧拉–麦克劳林求和公式如何在最重要的情形给出 \(s_n\) （或 \(r_n\) ）的任意阶渐近展开。

